% Generated draft: requires mathematical and source review.
\documentclass[11pt]{article}
\usepackage[margin=25mm]{geometry}
\usepackage{fontspec}
\setmainfont{Latin Modern Roman}
\usepackage{amsmath,amssymb,mathtools,unicode-math}
\setmathfont{Latin Modern Math}
\usepackage{enumitem,needspace,xurl,hyperref}
\setlength{\parindent}{0pt}
\setlength{\parskip}{5pt}
\setlength{\emergencystretch}{3em}
\newcommand{\N}{\mathbb{N}}
\newcommand{\Z}{\mathbb{Z}}
\newcommand{\Q}{\mathbb{Q}}
\newcommand{\R}{\mathbb{R}}
\newcommand{\C}{\mathbb{C}}
\newcommand{\F}{\mathbb{F}}
\newcommand{\A}{\mathbb{A}}
\newcommand{\PP}{\mathbb{P}}
\newcommand{\E}{\mathbb{E}}
\newcommand{\eps}{\varepsilon}
\newcommand{\dd}{\,\mathrm{d}}
\begin{document}
\section*{UnsolvedMath ID: 30001207}
\textbf{Source problem code:} catalogue entry
\subsection{Definitions and mathematical statement}
The supplied record asks for a construction and conditional geometric properties, but does not specify enough data to formulate a unique quantified conjecture.

Let \(b\) denote a Brun expansion in the sense intended by the source. The requested construction is an association
\[
 b\longmapsto\mathcal R(b),
\]
where \(\mathcal R(b)\) is a Rauzy fractal with associated subtiles \(\mathcal R_i(b)\), indexed by a set \(I\). The dimension, version of the Brun algorithm, admissible expansions, construction of the fractal, indexing of subtiles, and ambient space are not specified in the supplied record. This notation records the requested objects; it does not supply their missing definitions or assert their existence.

Once an ambient measured space \((E,\mu)\) and measurable subtiles have been specified, disjointness in measure means
\[
 \mu\bigl(\mathcal R_i(b)\cap\mathcal R_j(b)\bigr)=0
 \qquad(i,j\in I,\ i\ne j).
\]
The problem asks whether this property follows under the stated assumption of ``almost-everywhere exponential convergence to that fractal.'' The record does not identify the converging objects, their metric, the parameter space and measure underlying ``almost everywhere,'' or the dependence of exponential rates and constants on parameters. Consequently, that assumption cannot be made into a precise inequality from the supplied evidence.

The further task is to determine when the associated subtiles form a tiling. If the intended tiles are translates in a real vector space \(E\), one would need translation sets \(\Lambda_i\subseteq E\) and the indexed family
\[
 T_{i,t}=\mathcal R_i(b)+t,
 \qquad i\in I,\quad t\in\Lambda_i.
\]
For such a specified measurable family, coverage and disjointness up to measure zero would be expressed by
\[
 \mu\left(E\setminus\bigcup_{i\in I}\bigcup_{t\in\Lambda_i}T_{i,t}\right)=0,
 \qquad
 \mu(T_{i,t}\cap T_{j,s})=0
 \quad\text{for }(i,t)\ne(j,s).
\]
These formulas only explain one possible measure-theoretic meaning of tiling. The source does not specify translations, a lattice, the space to be tiled, or whether exact coverage, local finiteness, or additional boundary conditions are required. Exact coverage would instead require the displayed union to equal \(E\). Neither convention is selected here.

Thus the source-level task is to construct a Rauzy fractal from a Brun expansion, investigate measure-disjointness of its subtiles under the stated convergence assumption, and characterize its tiling properties. The missing definitions, hypotheses, and quantifiers require source clarification before this becomes a fully precise mathematical assertion.
\subsection{Short English statement}
Associate a Rauzy fractal and subtiles to a Brun expansion, then investigate whether the supplied almost-everywhere exponential convergence assumption implies measure-disjointness and when a tiling results. The algorithm, convergence hypothesis, measures, and tiling convention remain unspecified, so this is a qualified statement-only draft requiring mathematical review.
\subsection{Sources}
\begin{itemize}
\item Local UnsolvedMath registry, supplied source record: permanent ID \(30001207\), source problem code \texttt{catalogue entry}, title \emph{Rauzy Fractals from Multidimensional Continued Fractions}, collection \emph{Oberwolfach Reports Open Problems}. The record gives proposed year 2009 and names no individual proposer. Its statement is labeled \texttt{exact}; the supplied verification says it was retained byte-for-byte because it is internally readable. This is not independent source verification.
\item \emph{Mini-Workshop: The Pisot Conjecture --- From Substitution Dynamical Systems to Rauzy Fractals and Meyer Sets} (2009), \emph{Oberwolfach Reports}, pp.\ 725--766, DOI \texttt{10.4171/owr/2009/13}. Supplied source report and literature reference labeled ``original OWR report'': \url{https://doi.org/10.4171/owr/2009/13}. The record describes Oberwolfach Reports as documenting mathematical workshops and their open-problem discussions.
\item Supplied literature reference labeled ``Berthe survey with open question'': \url{https://www.irif.fr/~berthe/Articles/frac-arith-Leiden.pdf}. No full bibliographic details or pinpoint location are supplied.
\item Literature assessment dated exactly \texttt{2026-08-21}: the supplied curation labels the problem open and reports that the Brun-expansion Rauzy-fractal construction and tiling properties are explicitly left open, with later surveys continuing to describe the corresponding problem as open. This is supplied dated curation derived from the OWR record, not an independently verified literature assessment. The record explicitly requests independent verification of time-sensitive status and cited claims.
\end{itemize}
\end{document}
